\documentclass[11pt,a4paper]{article}

\usepackage[utf8]{inputenc}
\usepackage[T1]{fontenc}
\usepackage{amsmath,amssymb}
\usepackage{booktabs}
\usepackage{siunitx}
\usepackage{graphicx}
\usepackage{longtable}
\usepackage[margin=2.5cm]{geometry}
\usepackage[hidelinks]{hyperref}

\sisetup{detect-weight=true,detect-family=true}

\newcommand{\pu}{\ensuremath{\mathrm{p.u.}}}
\newcommand{\Mvar}{\ensuremath{\mathrm{Mvar}}}
\newcommand{\qss}{\textsc{qss}}
\newcommand{\rms}{\textsc{rms}}

\title{On the divergence between the quasi-steady-state and RMS
representations of a hierarchical reactive-power control cascade\\[2mm]
\large Candidate causes, experimental elimination, and the residual mechanism}
\author{}
\date{2026--07--25}

\begin{document}
\maketitle

\begin{abstract}
The quasi-steady-state (\qss) substitution premise underlying the two-layer
online feedback optimisation (OFO) cascade asserts that the reactive-power
plant may be represented by an algebraic power-flow solution rather than by
its electromechanical dynamics. Testing this premise against a PowerFactory
\rms{} model of a modified IEEE~39-bus system revealed a persistent
divergence between the two representations, concentrated in one distribution
underlay. This report enumerates the candidate causes, documents the
experiments used to eliminate them, and reports the residual mechanism. The
outcome is that the two plant models are equivalent to
\SI{0.106}{\Mvar} when the DER reactive-power droop is continuous, and that
the divergence is created by the \emph{dead zone} of the droop
characteristic, which admits more than one self-consistent equilibrium. Both
the \qss{} operating point and the point reached by the \rms{} trajectory are
shown to be valid equilibria of the same \rms{} differential-algebraic model.
\end{abstract}

\tableofcontents

%======================================================================
\section{Setting and notation}

\subsection{Plant and control hierarchy}

The test system is a modified IEEE~39-bus network (scenario
\texttt{wind\_replace}) partitioned into three transmission control zones,
each of which hosts high-voltage underlays denoted $\mathrm{DSO}_1 \dots
\mathrm{DSO}_4$. Layer~1 (transmission) dispatches AVR set-points, on-load
tap changers (OLTC), switched shunts and transmission-connected DER every
\SI{180}{\second}. Layer~2 (\SI{110}{\kilo\volt}) tracks the reactive-power
set-point issued at the transmission--distribution interface using local
OLTCs and DER every \SI{20}{\second}. Both layers are mixed-integer
quadratic programmes acting on cached sensitivities; neither controller has
access to the plant equations.

Two plant realisations are compared under an identical controller stack:

\begin{description}
  \item[\qss{} plant] a \texttt{pandapower} power flow with an anchored ZIP
        load model, solved to convergence at each dispatch instant
        (\texttt{run\_control=True}, i.e.\ the local droop loops are iterated
        to their fixed point);
  \item[\rms{} plant] a PowerFactory \rms{} time-domain model with WECC
        converter models (REGC/REEC), synchronous-machine AVRs and
        governors, integrated continuously between dispatch instants.
\end{description}

\subsection{The reactive-power droop law}

Each of the 44 DER parks realises a re-anchored, vertically shifted $Q(V)$
characteristic. With $S_n$ the park rating, $\sigma$ the droop slope,
$\delta$ the dead-zone half-width, $q^{\mathrm{set}}$ the OFO command and
$V^{\mathrm{a}}$ the anchor voltage,
%
\begin{equation}
  Q \;=\; \operatorname{clip}\!\Big(
    q^{\mathrm{set}} - \tfrac{S_n}{\sigma}\,
    \mathrm{dz}_{\delta}\!\left(V - V^{\mathrm{a}}\right),\;
    Q_{\min},\, Q_{\max}\Big),
  \label{eq:droop}
\end{equation}
%
where the dead-zone operator is
%
\begin{equation}
  \mathrm{dz}_{\delta}(x) =
  \begin{cases}
    x-\delta, & x > \delta,\\
    0, & |x| \le \delta,\\
    x+\delta, & x < -\delta .
  \end{cases}
  \label{eq:deadzone}
\end{equation}
%
At every dispatch instant both plants re-anchor $V^{\mathrm{a}} \leftarrow V$
and write the new offset $q^{\mathrm{set}}$, so that the droop contribution
is reset to zero and $Q = q^{\mathrm{set}}$ holds instantaneously; the droop
subsequently acts only on deviations from that operating point. Nominal
parameters are $\sigma = 0.06$, $\delta = 0.01\,\pu$.

In the \rms{} model, \eqref{eq:droop} is implemented as a DSL block in the
plant-control slot of the WECC composite,
%
\begin{equation}
  \dot{x}_1 = \frac{u - x_1}{T_f}, \qquad
  Q^{\mathrm{ext}} = \operatorname{clip}\!\big(
  q^{\mathrm{set}} - K_{\mathrm{droop}}\,
  \mathrm{dz}_{\delta}(x_1 - V^{\mathrm{a}}),\,
  q_{\min},\, q_{\max}\big),
  \label{eq:qvpre}
\end{equation}
%
with $K_{\mathrm{droop}} = 1/\sigma$ in per-unit of the park rating and
$T_f = \SI{0.02}{\second}$. The filter constant is three orders of magnitude
below the \SI{20}{\second} dispatch window, so at a dispatch instant
$x_1 = u$ to numerical precision and \eqref{eq:qvpre} reduces exactly to
\eqref{eq:droop}.

\subsection{Metric}

Unless stated otherwise the reported quantity is the root-mean-square
deviation between the \qss{} and \rms{} realisations of the twelve
transmission--distribution interface reactive-power flows, evaluated over a
common horizon of 15 dispatch intervals ($\SI{300}{\second}$, $n = 180$
samples). Using a fixed horizon is essential: the deviation grows with time,
so runs of different length are not comparable.

%======================================================================
\section{The observation}

Under closed-loop co-simulation with time-varying profiles, the two
representations do not track one another. The deviation is concentrated in
$\mathrm{DSO}_4$, whose DER reactive power diverges monotonically from
$\SI{2.9}{\Mvar}$ at $t=\SI{20}{\second}$ to $\SI{48}{\Mvar}$ after three
hours, after which it saturates. The remaining underlays exhibit a
transient deviation at the first profiled interval which subsequently
self-corrects. Because both plants are driven by one controller stack of
identical construction, and because the initial states agree to
$1.5\times10^{-5}\,\pu$, this behaviour requires explanation.

%======================================================================
\section{Candidate causes}

Four classes of explanation were considered. They are listed here in the
order in which they were tested, which is roughly the order of increasing
subtlety.

\begin{description}
  \item[\textbf{A. Command delivery.}] The \rms{} plant receives controller
        commands as discrete simulation events. If events are delivered
        late, partially, or not at all, the \rms{} plant executes a
        different input trajectory than the \qss{} plant, and no statement
        about plant equivalence can be made.
  \item[\textbf{B. Effective control law.}] The two implementations might
        realise different laws despite sharing the label ``$Q(V)$
        droop'': a different regulated voltage (remote, compensated, or
        filtered), proportional--integral rather than droop action,
        different gain normalisation, different limits, or a different load
        model. A related structural difference is the representation of
        primary frequency control.
  \item[\textbf{C. Numerical and initialisation effects.}] The \qss{}
        solution is obtained by a damped fixed-point iteration from a
        warm start. A different warm start, a different damping factor, or a
        different initial anchor could select a different numerical
        solution.
  \item[\textbf{D. Structural properties of the closed loop.}] The \rms{}
        endpoint might not be an equilibrium at all; more than one
        equilibrium might exist; or the two independent control loops might
        amplify an arbitrarily small plant difference.
\end{description}

%======================================================================
\section{Experimental elimination}

\subsection{A. Command delivery}

\paragraph{Event starvation.} PowerFactory admits newly created simulation
events into a running \rms{} calculation at a bounded rate. With profiles
enabled, approximately 177 events per dispatch interval were created, so the
backlog grew by roughly one interval per interval; events scheduled beyond
$t\approx\SI{300}{\second}$ of a \SI{600}{\second} run never executed. This
was established by row-exact forensics on the event folder and result file
of an affected run. The defect was removed by pre-creating a persistent event
pool before the calculation starts and by delivering exogenous profiles
through \texttt{ElmFile} playback rather than as events.

\paragraph{Re-anchor and offset delivery.} Each dispatch writes \emph{two}
parameter events per park: the offset $q^{\mathrm{set}}$ and the anchor
$V^{\mathrm{a}}$. Partial delivery of this pair would apply a new offset
against a stale anchor, injecting a spurious droop term
$K_{\mathrm{droop}}(V - V^{\mathrm{a}}_{\text{old}})$; at
$K_{\mathrm{droop}} = 16.7$ a stale anchor of $0.01\,\pu$ corresponds to a
reactive-power error of $17\%$ of park rating. The anchor signal is not
recorded in the standard monitor set, so this could not be excluded from
existing results.

A dedicated probe issued known dispatches and read the parameters back from
the DSL block (\texttt{c:qset}, \texttt{c:Vanchor}, \texttt{c:Kdroop},
\texttt{c:db}), then compared the delivered reactive power against
\eqref{eq:qvpre} evaluated with the parameters the simulator actually holds.
Results are given in Table~\ref{tab:delivery}. Both parameters are delivered
exactly, on all 44 parks, in every dispatch, under production
(\texttt{ElmFile}) and under the legacy event-based profile delivery. The
event pool showed no admission backlog
(\texttt{pending\_admission}~$=0$) with $5.2\times10^4$ objects present.

\begin{table}[htbp]
  \centering
  \caption{Delivery of the re-anchor pair, 44 parks $\times$ 6 dispatches.
  $r_Q$ is the residual of the assumed law evaluated with the parameters
  held by the simulator.}
  \label{tab:delivery}
  \begin{tabular}{lccc}
    \toprule
    Profile delivery & $\max|\Delta q^{\mathrm{set}}|$ [\pu] &
    $\max|\Delta V^{\mathrm{a}}|$ [\pu] & $\max |r_Q|$ [\Mvar] \\
    \midrule
    none                       & $0$                  & $0$ & $0.212$ \\
    \texttt{ElmFile} (production) & $1.4\times10^{-17}$ & $0$ & $0.043$ \\
    events (legacy, stress)    & $1.4\times10^{-17}$   & $0$ & $0.043$ \\
    \bottomrule
  \end{tabular}
\end{table}

The residual $r_Q \le \SI{0.043}{\Mvar}$ on parks rated several hundred
\si{\mega\volt\ampere} ($\approx 0.01\%$) is attributable to the difference
between the instantaneous terminal voltage used in the check and the
filtered state $x_1$ of \eqref{eq:qvpre} while the plant is still settling
within the advance. Its magnitude does not grow across dispatches.

\emph{Conclusion:} commands are delivered intact, and the \rms{} model
implements the assumed law. Class~A is eliminated.

\subsection{B. Effective control law}

Each candidate mismatch was checked against the implementation and, where
possible, against measurement. Table~\ref{tab:law} summarises.

\begin{table}[htbp]
  \centering
  \caption{Candidate law mismatches and their disposition.}
  \label{tab:law}
  \begin{tabular}{p{0.32\linewidth}p{0.50\linewidth}c}
    \toprule
    Candidate & Finding & Status \\
    \midrule
    Remote or pilot-node voltage &
    \rms{} reads the park's own terminal (\texttt{bus1.cterm}); \qss{} reads
    the same bus; parity $1.5\times10^{-5}$ & excluded \\
    Compensated voltage ($X_c$) &
    $v_{\mathrm{eff}} = x_1 - V^{\mathrm{a}}$ contains no current term &
    excluded \\
    Filtered voltage &
    $T_f = \SI{0.02}{\second} \ll \SI{20}{\second}$; no steady-state lag &
    excluded \\
    PI regulator instead of droop &
    the DSL block has no integrator on the control path; both sides permit
    $V \neq V_{\mathrm{ref}}$ & excluded \\
    Gain normalisation &
    $K_{\mathrm{droop}} = 1/\sigma$ equals $S_n/\sigma$ in per-unit &
    excluded \\
    Dead zone &
    identical parameter on both sides & \emph{shared, see \S\ref{sec:mech}} \\
    Reactive limits &
    mirrored from one capability function; diagnostic runs use a
    $\pm1.0\,\pu$ box so limits do not bind & excluded \\
    Load model &
    anchored ZIP on both sides & excluded \\
    Discrete states (taps, shunts) &
    zero tap divergences in the converged configuration & excluded \\
    Primary frequency control &
    \qss{} uses a distributed slack ($\Delta P \propto S_n$, no frequency
    state); \rms{} uses non-uniform governors & \textbf{structural,
    bounded} \\
    \bottomrule
  \end{tabular}
\end{table}

The frequency representation is a genuine structural difference and cannot
be excluded by inspection. It is, however, bounded by the open-loop
experiment of \S\ref{sec:uy}: with a continuous droop the two plants agree
to \SI{0.106}{\Mvar} under identical inputs, which bounds the aggregate
contribution of all remaining structural differences, frequency included.

\subsection{C. Numerical and initialisation effects}

Three experiments were carried out on the \qss{} side alone, and are
therefore free of any \rms{} interaction.

\paragraph{Warm-start dependence.} The closed-loop iteration was started
from three initial conditions: the analytical linear warm start, zero
reactive power, and the reactive power at which the \rms{} model settles.
All three converge to the same point
($\{26.7,\,13.3,\,7.5,\,22.8\}\,\Mvar$ per underlay), including the start
placed exactly at the \rms{} solution. The \qss{} equilibrium is therefore
independent of the warm start; the warm start accelerates convergence but
does not select the solution.

\paragraph{Damping dependence.} The local loop implements
$Q \leftarrow Q + \lambda\,(Q^{\star}(V) - Q)$, whose fixed points are
independent of $\lambda$. Sweeping $\lambda \in \{0.5, 0.1, 0.02\}$ from
both starting points confirms this: all convergent cases reach the same
point (residual $\approx \SI{0.09}{\Mvar}$). At $\lambda = 0.5$ the
iteration fails to converge and exhibits a limit cycle of approximately
\SI{17}{\Mvar} peak-to-peak with identical amplitude from both starts,
i.e.\ one cycle rather than a second attractor.

\paragraph{Initial anchor mismatch.} Prior to the first dispatch the \qss{}
droop falls back to a nominal anchor of $1.03\,\pu$, whereas the \rms{}
model is anchored at the local terminal voltage ($\approx 1.02\,\pu$).
A configuration option was implemented to anchor both sides at the local
voltage at initialisation. The resulting change is negligible: the
$\mathrm{DSO}_4$ deviation after \SI{600}{\second} is \SI{-18.3}{\Mvar}
without and \SI{-19.5}{\Mvar} with the correction.

\emph{Conclusion:} class~C is eliminated. The \qss{} solution is not an
artefact of warm start, damping, or initial anchoring.

\subsection{D.1 Stationarity of the RMS endpoint}
\label{sec:hold}

A simulation end time does not prove an equilibrium, and the system's
electromechanical modes ring for \SIrange{13}{22}{\second} against a
\SI{20}{\second} dispatch interval. The following test was therefore
carried out. The \qss{} plant was advanced to a profiled operating point and
its \emph{closed-loop} problem converged there, so that the droop is at its
fixed point. PowerFactory was synchronised to that state, initialised, and
then advanced for \SI{300}{\second} with no dispatch, no profile step and no
disturbance of any kind. Drift was monitored on bus voltages, park reactive
power, machine reactive power and machine speed.

The worst bus-voltage drift over \SI{300}{\second} is
$6.17\times10^{-11}\,\pu$, and it is flat: $6.11\times10^{-11}$ is already
reached after the first \SI{20}{\second} chunk. Reactive powers and machine
speeds do not move within printing precision. This test was run at the
nominal dead zone $\delta = 0.01\,\pu$, i.e.\ in the configuration in which
the closed loop diverges.

\emph{Conclusion:} the \qss{} operating point \emph{is} an equilibrium of
the \rms{} differential-algebraic model. Class~D.1 is eliminated.

%======================================================================
\section{The residual mechanism}
\label{sec:mech}

\subsection{Uniqueness of the continuous droop}

For a continuous droop the equilibrium is unique. Writing the network
response as $V = V_0 + S_{VQ} Q$ with $S_{VQ} = \partial V/\partial Q > 0$
and the control law without dead zone as
$Q = q^{\mathrm{set}} - R\,(V - V^{\mathrm{a}})$, $R = S_n/\sigma > 0$,
elimination gives
%
\begin{equation}
  Q\,\big(1 + R\,S_{VQ}\big)
  \;=\;
  q^{\mathrm{set}} - R\,\big(V_0 - V^{\mathrm{a}}\big),
  \label{eq:unique}
\end{equation}
%
and since $R S_{VQ} > 0$ the coefficient never vanishes. A steeper droop
increases the loop gain $R S_{VQ}$ and may render the fixed-point
\emph{iteration} oscillatory, but it does not create a second solution.

\subsection{Effect of the dead zone}

The dead zone destroys this argument, and does so even though it is
implemented identically in both models. Inside the dead zone the control law
reduces to $Q = q^{\mathrm{set}}$, independent of $V$: the park becomes a
constant reactive-power source and \eqref{eq:unique} no longer constrains
the voltage. The equilibrium condition is thus piecewise, and near an edge
$|V - V^{\mathrm{a}}| \approx \delta$ the in-band and out-of-band branches
can both be self-consistent. Which branch a park occupies is itself part of
the solution, and with 44 parks the number of admissible branch combinations
is large.

This is the central point, and it differs from the usual framing: the
divergence does not arise because the two models contain \emph{different}
dead zones, but because they contain the \emph{same} one. A correctly
implemented, shared dead zone is by itself sufficient to admit more than one
equilibrium.

\subsection{Which excursion actually reaches the dead-zone edge}
\label{sec:excursion}

The argument of the preceding subsection is necessary but not sufficient: it
establishes that multiple branches \emph{can} exist, not that the system
visits them. A screening of the \qss{} plant alone settles this and, in
doing so, corrects a plausible but incorrect reading of the mechanism.

Running the static closed loop over a three-hour profiled horizon and
recording, for every dispatch interval and park, the excursion from the
anchor set at that dispatch,
$\mathrm{exc}(i,k) = |V_i(t_{k+1}) - V^{\mathrm{a}}_i(t_k)|$, gives a
maximum of $0.00828\,\pu$ over $4.7\times10^4$ park-intervals, with the
$99$th percentile at $0.00178\,\pu$. At the nominal dead zone
$\delta = 0.01\,\pu$ the \qss{} droop therefore \emph{never} engages. Since a
dormant droop implies $Q \equiv q^{\mathrm{set}}$ and hence identical
injections, the divergence cannot be explained by the two plants selecting
different branches of the same characteristic.

The resolution is that the \emph{dynamic} excursion of the \rms{} model is
substantially larger than the quasi-static displacement, because the
integration passes through the electromechanical transient rather than
jumping to the algebraic solution. Evaluating the same statistic on the
harvested \rms{} terminal voltages gives the values in
Table~\ref{tab:excursion}.

\begin{table}[htbp]
  \centering
  \caption{Maximum per-interval voltage excursion in the \rms{} model,
  compared with the $\delta$-independent quasi-static value of
  $0.00828\,\pu$, and the resulting plant residual.}
  \label{tab:excursion}
  \begin{tabular}{lcccc}
    \toprule
    $\delta$ [\pu] & \rms{} max excursion [\pu] & ratio to $\delta$ &
    crossing rate & plant residual [\Mvar]\\
    \midrule
    $0$      & $0.00227$ & ---    & ---      & $0.362$\\
    $0.0025$ & $0.00452$ & $1.81$ & $20.9\%$ & $0.289$\\
    $0.005$  & $0.00695$ & $1.39$ & $12.7\%$ & $0.288$\\
    $0.0075$ & $0.00968$ & $1.29$ & $10.9\%$ & $0.592$\\
    $0.01$   & $0.01256$ & $1.26$ & $8.8\%$  & $0.690$\\
    $0.015$  & $0.01912$ & $1.27$ & $7.0\%$  & $8.352$\\
    $0.02$   & $0.02384$ & $1.19$ & $5.3\%$  & $8.851$\\
    \bottomrule
  \end{tabular}
\end{table}

Two regularities emerge. First, the \rms{} excursion is essentially
\emph{set} by the dead zone, $\mathrm{exc}_{\max} \approx 1.2$--$1.3\,\delta$
for $\delta \ge 0.005\,\pu$: the voltage drifts unopposed until it clears the
dead zone, at which point the droop engages and arrests it. The dead-zone
half-width therefore prescribes the amplitude of the dynamic voltage swing.
Second, the quasi-static displacement is $\delta$-independent at
$0.00828\,\pu$, because the algebraic model has no transient to damp.

At $\delta = 0.01\,\pu$ the droop consequently engages in $8.8\%$ of
\rms{} park-intervals and in none of the \qss{} ones. The two plants then
apply different reactive power in response to identical commands, which is
the immediate origin of the residual.

Within the reference operating condition these figures suggested a screening
criterion requiring no \rms{} simulation, namely $\delta \lesssim
\max_{i,k}|V_i(t_{k+1}) - V^{\mathrm{a}}_i(t_k)|_{\qss}$, whose right-hand
side is obtainable from a static run in minutes and evaluates to
$0.0083\,\pu$ here.

\paragraph{This criterion was tested out of sample and failed.} A
quasi-static screening across a season\,$\times$\,time-of-day grid showed the
right-hand side to vary by a factor of $2.52$ over the year
($0.00812$--$0.02051\,\pu$), so the criterion cannot be a fixed number. The
prediction was therefore registered in the sharper form: at
2016-07-15 03:00, where the quasi-static excursion is $0.02051\,\pu$ and the
\qss{} droop itself engages in $1.48\%$ of park-intervals (against $0\%$ at
the reference), the reduced \qss/\rms{} asymmetry should give a residual
\emph{below} the reference \SI{0.69}{\Mvar}. The measured residual is
\SI{8.28}{\Mvar}, a factor of twelve in the opposite direction. The criterion
is consequently withdrawn.

The same measurement also excludes the simpler reading that the residual
follows the \rms{} swing amplitude: at the summer condition the \rms{}
excursion is \emph{smaller} than at the reference ($0.01032$ versus
$0.01256\,\pu$) while the residual is twelve times larger. The sign of the
asymmetry even reverses --- at the reference the \rms{} moves further, at the
summer condition the quasi-static model does, its algebraic step overshooting
what the \rms{} attains dynamically within one interval.

A further hypothesis, that the residual is ordered by the \emph{difference}
between the per-interval displacements of the two representations, was also
tested out of sample and rejected: at 2016-01-15 03:00 that difference takes
its smallest measured value ($0.00199\,\pu$, against $0.0043$ at the
reference) while the residual is \SI{4.90}{\Mvar}, seven times the reference
value.

\paragraph{The quasi-static excursion orders both axes.} The analyses above
assumed the quasi-static displacement to be independent of $\delta$. Direct
measurement refutes that assumption: it varies by a factor of $8.5$ across
the dead-zone range, because a wider dead zone withholds regulation and
allows the operating point to travel further within a dispatch interval.
Recomputing with the measured values places every result obtained so far ---
seven dead-zone settings at one operating condition, and three operating
conditions at one dead-zone setting --- on a single monotone curve,
Table~\ref{tab:unified}.

\begin{table}[htbp]
  \centering
  \caption{Plant residual ordered by the quasi-static per-interval voltage
  excursion. Rows in bold are different \emph{operating conditions} at the
  nominal $\delta = 0.01\,\pu$; the remainder are dead-zone settings at the
  reference condition.}
  \label{tab:unified}
  \begin{tabular}{ccl}
    \toprule
    \qss{} excursion [\pu] & plant residual [\Mvar] & source \\
    \midrule
    $0.00255$ & $0.362$ & $\delta = 0$ \\
    $0.00463$ & $0.289$ & $\delta = 0.0025$ \\
    $0.00629$ & $0.288$ & $\delta = 0.005$ \\
    $0.00798$ & $0.592$ & $\delta = 0.0075$ \\
    $0.00828$ & $0.690$ & $\delta = 0.01$ (reference) \\
    $\mathbf{0.01573}$ & $\mathbf{4.903}$ & \textbf{2016-01-15 03:00} \\
    $0.01746$ & $8.352$ & $\delta = 0.015$ \\
    $\mathbf{0.02051}$ & $\mathbf{8.282}$ & \textbf{2016-07-15 03:00} \\
    $0.02161$ & $8.851$ & $\delta = 0.02$ \\
    \bottomrule
  \end{tabular}
\end{table}

The curve is a low plateau of $0.3$--$\SI{0.7}{\Mvar}$ below roughly
$0.009\,\pu$, a steep transition, and a high plateau of $8$--$\SI{9}{\Mvar}$
above roughly $0.017\,\pu$. The single departure from monotonicity is at the
lower end, where $\delta = 0$ gives a slightly larger residual than
$\delta = 0.0025$; this is the chattering regime discussed in
\S\ref{sec:design}, in which the undamped droop responds to arbitrarily small
deviations.

The operating-condition rows constitute an out-of-sample test, since the
curve is determined entirely by the dead-zone sweep at the reference
condition. Both are placed correctly: $0.02051\,\pu$ falls on the upper
plateau (measured \SI{8.28}{\Mvar}) and $0.01573\,\pu$ within the transition
(measured \SI{4.90}{\Mvar}). The screening criterion is therefore restated as
an absolute condition on the excursion itself,
%
\begin{equation}
  \max_{i,k}\,\big|V_i(t_{k+1}) - V^{\mathrm{a}}_i(t_k)\big|_{\qss}
  \;\lesssim\; 0.009\,\pu ,
  \label{eq:criterion}
\end{equation}
%
rather than as a condition relating $\delta$ to that excursion, which was the
form rejected above. The left-hand side is obtainable from a static run in
minutes and requires no \rms{} simulation; both the dead-zone width and the
operating condition enter through it.

The numerical thresholds in \eqref{eq:criterion} are calibrated on one
network and one profile year, and the steepness of the transition means that
excursions between $0.009$ and $0.017\,\pu$ admit no quantitative prediction
of the residual --- only the statement that it is neither small nor
saturated.

\subsection{Positive demonstration of multiplicity}

Two facts established above combine into a direct demonstration:

\begin{enumerate}
  \item the \qss{} operating point is an equilibrium of the \rms{} model
        (\S\ref{sec:hold}, drift $6.17\times10^{-11}\,\pu$ over
        \SI{300}{\second});
  \item the \rms{} model, integrated along the co-simulation trajectory,
        settles at a \emph{different} point.
\end{enumerate}

Both points are solutions of the same model. The system therefore possesses
at least two equilibria, and the trajectory selects between them. This is a
constructive demonstration rather than an inference by elimination.

%======================================================================
\section{Quantitative results}
\label{sec:uy}

\subsection{Plant equivalence under identical inputs}

To separate plant equivalence from control-loop behaviour, the actuator and
profile trajectory recorded from the \qss{} run was replayed verbatim into
the \rms{} plant, so that both plants receive identical inputs and no
controller can react. The resulting residual is a property of the plant
models alone. It was measured as a function of the dead-zone half-width.

\subsection{Closed-loop deviation and amplification}

The same sweep was repeated in closed loop, i.e.\ with two independent
controller stacks. The ratio of the closed-loop deviation to the plant
residual isolates the contribution of control-loop interaction.
Table~\ref{tab:map} collects both, together with the droop-gain sweep.

\begin{table}[htbp]
  \centering
  \caption{Validity map under the $\pm1.0\,\pu$ diagnostic capability
  override. Interface-Q RMSE over 15 dispatch intervals.}
  \label{tab:map}
  \begin{tabular}{lccc}
    \toprule
    Configuration & Plant residual [\Mvar] & Closed loop [\Mvar] &
    Amplification \\
    \midrule
    \multicolumn{4}{l}{\emph{Dead-zone sweep} ($\sigma = 0.06$)}\\
    $\delta = 0$        & $0.106$ & $0.12$ & $1.1$ \\
    $\delta = 0.005$    & $0.246$ & $0.65$ & $2.6$ \\
    $\delta = 0.01$     & $0.710$ & $2.11$ & $3.0$ \\
    $\delta = 0.02$     & $8.906$ & $4.33$ & $0.5$ \\
    \midrule
    \multicolumn{4}{l}{\emph{Droop-gain sweep} ($\delta = 0.01$)}\\
    $\sigma = 0.12$ (half gain)   & $0.710$ & $1.26$ & $1.8$ \\
    $\sigma = 0.06$ (nominal)     & $0.710$ & $2.11$ & $3.0$ \\
    $\sigma = 0.03$ (double gain) & $0.710$ & $9.04$ & $12.7$ \\
    \bottomrule
  \end{tabular}
\end{table}

Three observations follow.

\begin{enumerate}
  \item \textbf{With a continuous droop the two plant models are
        equivalent.} At $\delta = 0$ the residual under identical inputs is
        \SI{0.106}{\Mvar}, which is the numerical floor of the comparison.
        Any structural mismatch --- including the frequency representation
        --- is bounded by this figure.
  \item \textbf{The dead zone is the dominant driver.} The plant residual
        grows by a factor of approximately $84$ across the range
        $\delta \in [0, 0.02]$. This is a plant-level effect, present with
        no controller in the loop.
  \item \textbf{The control-loop contribution is secondary and changes
        sign.} For narrow dead zones the two independent loops amplify the
        plant residual (up to $3\times$, and $12.7\times$ at double droop
        gain). At $\delta = 0.02$ the amplification falls below unity: the
        wide dead zone disables local voltage regulation, so under identical
        inputs nothing holds the two trajectories together, whereas in
        closed loop the OFO controllers observe the drift and correct it.
        The curves cross at $\delta \approx 0.0125$.
\end{enumerate}

\subsection{Consequence for the closed-loop cascade}

Removing the dead zone restores agreement in full closed-loop operation. A
one-hour co-simulation with $\delta = 0$ yields an interface-Q RMSE of
\SI{1.03}{\Mvar} and a zone-voltage RMSE of $8.5\times10^{-4}\,\pu$, with no
unsettled dispatch interval out of 2160 interface-power and 540 voltage
samples, and no divergence of discrete states.

\subsection{Results under physical capability limits}
\label{sec:vde}

The results of Table~\ref{tab:map} were obtained with a diagnostic
$\pm1.0\,\pu$ reactive-capability box, which removes limit binding and
therefore isolates the droop mechanism. The sweep was repeated under the
physical VDE-AR-N-4120 and STATCOM operating diagrams on a finer dead-zone
grid; under those limits several parks saturate and the distribution coupler
OLTCs become active. Table~\ref{tab:vde} and Figure~\ref{fig:vde} report the
outcome.

\begin{table}[htbp]
  \centering
  \caption{Validity map under physical VDE capability, compared with the
  $\pm1.0\,\pu$ diagnostic override. Interface-Q RMSE over 15 dispatch
  intervals, $\sigma = 0.06$.}
  \label{tab:vde}
  \begin{tabular}{lccc@{\hskip 8mm}ccc}
    \toprule
    & \multicolumn{3}{c}{physical VDE}
    & \multicolumn{3}{c}{$\pm1.0\,\pu$ diagnostic}\\
    \cmidrule(lr){2-4}\cmidrule(lr){5-7}
    $\delta$ [\pu] & floor & closed & amplif.
                   & floor & closed & amplif.\\
    \midrule
    $0$      & $0.362$ & $0.599$ & $1.65$ & $0.106$ & $0.119$ & $1.12$\\
    $0.0025$ & $0.289$ & $0.493$ & $1.71$ & ---     & ---     & ---   \\
    $0.005$  & $0.288$ & $0.600$ & $2.09$ & $0.246$ & $0.645$ & $2.63$\\
    $0.0075$ & $0.592$ & $2.053$ & $3.47$ & ---     & ---     & ---   \\
    $0.01$   & $0.690$ & $2.099$ & $3.04$ & $0.710$ & $2.113$ & $2.98$\\
    $0.015$  & $8.352$ & $5.445$ & $0.65$ & ---     & ---     & ---   \\
    $0.02$   & $8.851$ & $4.206$ & $0.48$ & $8.906$ & $4.325$ & $0.49$\\
    \bottomrule
  \end{tabular}
\end{table}

\begin{figure}[htbp]
  \centering
  \includegraphics[width=\linewidth]{../results/p2_validity_map/vde_validity_map.pdf}
  \caption{Plant equivalence, closed-loop deviation and amplification as
  functions of the dead-zone half-width, under physical and diagnostic
  capability models. Logarithmic ordinate in (a) and (b).}
  \label{fig:vde}
\end{figure}

Four conclusions follow, of which the third was not anticipated.

\begin{enumerate}
  \item \textbf{Capability limits raise the plant residual only where the
        dead zone is inactive.} At $\delta = 0$ the residual increases from
        \SI{0.106}{\Mvar} to \SI{0.362}{\Mvar} when physical limits are
        restored, a factor of $3.4$: saturation pins the reactive power of
        individual parks and thereby removes part of the local voltage
        regulation. At $\delta \ge 0.01$ the two capability models give
        indistinguishable residuals ($0.690$ versus $0.710$, and $8.851$
        versus \SI{8.906}{\Mvar}). Once the dead zone is wide enough, it
        dominates completely and the capability model is irrelevant.
  \item \textbf{The degradation is a threshold, not a gradual trend.}
        Between $\delta = 0.01$ and $\delta = 0.015$ the plant residual
        rises from \SI{0.690}{\Mvar} to \SI{8.352}{\Mvar}, a factor of
        $12$ in a single step, while the four points below $\delta = 0.01$
        remain within a factor of $2.4$ of one another. The finer grid
        introduced for this sweep was necessary to resolve this; the
        coarser grid of Table~\ref{tab:map} suggested a smooth power law.
  \item \textbf{The nominal dead zone lies immediately below the
        threshold.} The design value $\delta = 0.01\,\pu$ yields a plant
        residual of \SI{0.690}{\Mvar}, whereas $\delta = 0.015\,\pu$ yields
        \SI{8.35}{\Mvar}. The quasi-steady-state substitution is therefore
        admissible at the nominal parameterisation, but without margin: a
        50\,\% increase of the dead zone invalidates it.
  \item \textbf{The amplification structure is independent of the
        capability model.} Both capability regimes peak at
        $3.0$--$3.5$ near $\delta \approx 0.0075$--$0.01$ and fall below
        unity between $\delta = 0.01$ and $\delta = 0.015$. The
        control-loop mechanism of \S\ref{sec:mech} is thus a property of
        the dead-zone characteristic and not an artefact of the diagnostic
        capability box under which it was first observed.
\end{enumerate}

%======================================================================
\section{The dead zone as a design parameter}
\label{sec:design}

The preceding sections treat the dead-zone half-width as a property of the
plant to be respected, not as a quantity to be chosen. It must be
emphasised that selecting $\delta$ so as to improve the agreement between
the two models would invert the purpose of the study: the dead zone is a
grid-code parameter, and if the installed parks operate at
$\delta = 0.01\,\pu$ then the model must do likewise. The result of
\S\ref{sec:vde} is a condition on the validity of the method, not a
recommendation to modify the plant.

If, however, $\delta$ is genuinely a design degree of freedom, the relevant
criteria are control performance and actuator effort rather than model
agreement. Table~\ref{tab:design} evaluates both on the \rms{} plant of the
physical-capability sweep.

\begin{table}[htbp]
  \centering
  \caption{Dead zone as a design choice: control performance and actuator
  effort, physical VDE capability, \rms{} plant, 15 dispatch intervals.
  Tracking error is $|Q^{\mathrm{act}} - Q^{\mathrm{set}}|$ over the twelve
  interfaces; travel is the summed inter-dispatch change of DER reactive
  power per underlay.}
  \label{tab:design}
  \begin{tabular}{lcccc}
    \toprule
    $\delta$ [\pu] & mean $|e|$ [\Mvar] & max $|e|$ [\Mvar]
                   & DER-$Q$ travel [\Mvar] & tap switches \\
    \midrule
    $0$        & $2.846$ & $6.337$  & $100.8$ & $0$ \\
    $0.0025$   & $1.242$ & $6.187$  & $104.7$ & $0$ \\
    $0.005$    & $\mathbf{1.104}$ & $6.250$ & $\mathbf{100.1}$ & $0$ \\
    $0.0075$   & $1.496$ & $6.815$  & $110.6$ & $0$ \\
    $0.01$     & $1.652$ & $10.179$ & $126.4$ & $1$ \\
    $0.015$    & $3.013$ & $16.826$ & $156.2$ & $1$ \\
    $0.02$     & $2.912$ & $16.828$ & $145.2$ & $1$ \\
    \bottomrule
  \end{tabular}
\end{table}

Two observations are relevant.

First, the tracking error is minimised at $\delta \approx 0.005\,\pu$, where
it is $33\%$ below the value at the nominal $\delta = 0.01\,\pu$, at
essentially unchanged converter travel and without tap operation. The plant
residual of \S\ref{sec:vde} is minimised at the same width. The two criteria
therefore agree.

Second, and less intuitively, $\delta = 0$ is a poor operating choice
despite being the configuration in which the multiplicity of \S\ref{sec:mech}
disappears: the tracking error degrades to \SI{2.85}{\Mvar}, a factor of
$2.6$ worse than at $\delta = 0.005\,\pu$. Without a dead zone every park
applies the full droop gain $K_{\mathrm{droop}} = 16.7$ to arbitrarily small
voltage deviations, including those arising from the mutual interaction of
the 44 parks, so the local loops chatter against one another. A narrow dead
zone suppresses this without materially reducing the regulating action. The
configuration that is most convenient for model comparison is thus not the
configuration that controls best, and the two should not be conflated.

The co-simulation is deterministic --- measurement noise is disabled,
profiles are fixed and the mixed-integer programmes are solved to optimality
--- so the differences in Table~\ref{tab:design} are exact for this
configuration rather than samples requiring an error estimate. They
nevertheless correspond to a single operating point and a \SI{300}{\second}
window; a recommendation would require the comparison to be repeated across
operating points and over a full profile day, since a narrower dead zone
implies response to smaller excursions and hence potentially greater
converter duty over long horizons.

%======================================================================
\section{Consequences for the premises}

The premise under test, $y_{\rms}(\infty; u) = y_{\qss}(u)$, must be
qualified rather than accepted or rejected outright.

\begin{itemize}
  \item It holds to \SI{0.106}{\Mvar} when the DER characteristic is
        continuous, and to \SI{0.362}{\Mvar} when physical capability
        limits are additionally imposed. In that regime the
        quasi-steady-state substitution is justified, and the closed-loop
        cascade reproduces the dynamic result.
  \item Its validity is bounded by a \emph{threshold} in the dead-zone
        half-width rather than degrading smoothly: at the reference
        operating condition the plant residual rises by a factor of $12$
        between $\delta = 0.01$ and $\delta = 0.015\,\pu$.
  \item \textbf{That threshold is not a property of the system alone.} At a
        different operating condition (2016-07-15 03:00) the residual is
        \SI{8.28}{\Mvar} already at $\delta = 0.01\,\pu$, i.e.\ the
        parameterisation that is admissible at the reference condition is
        not admissible there. Any statement of the form ``the premise holds
        for $\delta \le \delta^{\ast}$'' therefore requires the operating
        condition to be quantified as well; a single-condition study
        overstates the validity region.
  \item The nominal design value $\delta = 0.01\,\pu$ lies immediately
        below the threshold \emph{at the reference condition}, and above it
        at the summer condition tested. The premise is thus conditional on
        the operating point, which should be stated explicitly rather than
        left implicit.
  \item A separate premise --- that two independently controlled
        realisations of the same plant remain comparable --- fails
        independently, with a magnitude proportional to the droop loop gain
        (up to $12.7\times$ at double gain) and with a structure that is
        invariant to the capability model.
\end{itemize}

The practical implication is that the substitution premise should be stated
with a validity condition on the actuator characteristic, rather than as an
unconditional property of the timescale separation.

%======================================================================
\section{Threats to validity}

\begin{itemize}
  \item All results correspond to a single profile day and one operating
        point; the dead-zone effect is expected to depend on where the
        operating point sits relative to the anchor.
  \item The comparison horizon is \SI{300}{\second} for the sweep and
        \SI{3600}{\second} for the confirmation run. The deviation grows
        with time before saturating, so absolute figures are
        horizon-dependent; the reported ratios are not.
  \item The demonstration of multiplicity establishes that at least two
        equilibria exist and that the trajectory selects between them. It
        does not characterise the basins of attraction, nor does it
        enumerate the equilibria.
  \item Stationarity was verified on bus voltages, park and machine
        reactive power, and machine speed. Limiter states and tap timers
        were not monitored individually; they are inactive in the
        configuration tested.
  \item The frequency representation differs structurally between the two
        plants. Its contribution is bounded by \SI{0.106}{\Mvar} in the
        present study, but it would not be negligible in an
        active-power-dominated investigation.
\end{itemize}

\end{document}
